\textbf{Phần 1. Trắc nghiệm lựa chọn}
\setcounterpageref{ex}{0}

\begin{ex} %[3D2V24]
Năm nay con 7 tuổi, tuổi của bố gấp 5 lần tuổi của con. Hỏi năm nay bố bao nhiêu tuổi?
\choice
{$30$ tuổi}
{\True $35$ tuổi}
{$40$ tuổi}
{$42$ tuổi}
\loigiai{
Tuổi bố năm nay là: $7 \times 5 = 35$ (tuổi).
\textbf{Đáp án của câu hỏi là B}
}
\end{ex}

\begin{ex} %[3D2V24]
Một sợi dây đỏ dài $14\text{ cm}$. Sợi dây xanh dài gấp 4 lần sợi dây đỏ. Cả hai sợi dây dài tất cả là
\choice
{$56\text{ cm}$}
{$64\text{ cm}$}
{$72\text{ cm}$}
{\True $70\text{ cm}$}
\loigiai{
Sợi dây xanh dài: $14 \times 4 = 56\text{ (cm)}$.\
Cả hai sợi dây dài: $14 + 56 = 70\text{ (cm)}$.
\textbf{Đáp án của câu hỏi là D}
}
\end{ex}

\begin{ex} %[3D2V24]
Số lớn gấp 6 lần số bé. Biết số bé là 12, hiệu giữa hai số là
\choice
{\True $60$}
{$72$}
{$54$}
{$48$}
\loigiai{
Số lớn là: $12 \times 6 = 72$. Hiệu hai số là: $72 - 12 = 60$.
\textbf{Đáp án của câu hỏi là A}
}
\end{ex}

\begin{ex} %[3D2V24]
Muốn gấp một số lên một số lần, ta làm như sau:
\choice
{Lấy số đó cộng với số lần}
{Lấy số đó chia cho số lần}
{\True Lấy số đó nhân với số lần}
{Lấy số đó trừ đi số lần}
\loigiai{
Quy tắc toán học: Muốn gấp một số lên nhiều lần, ta lấy số đó nhân với số lần.
\textbf{Đáp án của câu hỏi là C}
}
\end{ex}

\begin{ex} %[3D2V24]
Thùng thứ nhất đựng $16\ell$ nước mắm. Thùng thứ hai đựng số nước mắm gấp 3 lần thùng thứ nhất. Thùng thứ hai đựng nhiều hơn thùng thứ nhất bao nhiêu lít nước mắm?
\choice
{$48\ell$}
{\True $32\ell$}
{$24\ell$}
{$36\ell$}
\loigiai{
Thùng thứ hai đựng: $16 \times 3 = 48\text{ (\ell)}$. Thùng thứ hai nhiều hơn: $48 - 16 = 32\text{ (\ell)}$.
\textbf{Đáp án của câu hỏi là B}
}
\end{ex}

\begin{ex} %[3D2V24]
Một số gấp lên 3 lần rồi cộng thêm 14 thì được 50. Số đó là
\choice
{$10$}
{$14$}
{$16$}
{\True $12$}
\loigiai{
Số đó gấp lên 3 lần là: $50 - 14 = 36$. Số ban đầu là: $36 : 3 = 12$.
\textbf{Đáp án của câu hỏi là D}
}
\end{ex}

\begin{ex} %[3D2V24]
Đoạn thẳng $AB$ dài $9\text{ cm}$. Đoạn thẳng $CD$ dài gấp 5 lần đoạn thẳng $AB$. Đoạn thẳng $MN$ dài ngắn hơn đoạn thẳng $CD$ là $8\text{ cm}$. Độ dài đoạn thẳng $MN$ là
\choice
{\True $37\text{ cm}$}
{$45\text{ cm}$}
{$35\text{ cm}$}
{$42\text{ cm}$}
\loigiai{
Đoạn thẳng $CD$ dài: $9 \times 5 = 45\text{ (cm)}$. Đoạn thẳng $MN$ dài: $45 - 8 = 37\text{ (cm)}$.
\textbf{Đáp án của câu hỏi là A}
}
\end{ex}

\begin{ex} %[3D2V24]
Một cửa hàng buổi sáng bán được 18 gói kẹo, buổi chiều bán được số gói kẹo gấp đôi buổi sáng. Hỏi cả ngày cửa hàng bán được bao nhiêu gói kẹo?
\choice
{$36$ gói kẹo}
{$48$ gói kẹo}
{\True $54$ gói kẹo}
{$60$ gói kẹo}
\loigiai{
Buổi chiều bán được: $18 \times 2 = 36$ gói. Cả hai buổi bán được: $18 + 36 = 54$ gói kẹo.
\textbf{Đáp án của câu hỏi là C}
}
\end{ex}

\begin{ex} %[3D2V24]
Một số gấp lên 4 lần thì bằng tích của 16 và 2. Số đó là
\choice
{$6$}
{\True $8$}
{$10$}
{$12$}
\loigiai{
Tích là: $16 \times 2 = 32$. Số cần tìm là: $32 : 4 = 8$.
\textbf{Đáp án của câu hỏi là B}
}
\end{ex}

\begin{ex} %[3D2V24]
Có ba bao thóc. Bao thứ nhất nặng $20\text{ kg}$. Bao thứ hai nặng gấp đôi bao thứ nhất. Bao thứ ba nặng bằng một nửa bao thứ hai. Tổng cân nặng của cả ba bao là
\choice
{$60\text{ kg}$}
{$70\text{ kg}$}
{$90\text{ kg}$}
{\True $80\text{ kg}$}
\loigiai{
Bao hai nặng: $20 \times 2 = 40\text{ kg}$. Bao ba nặng: $40 : 2 = 20\text{ kg}$. Tổng ba bao: $20 + 40 + 20 = 80\text{ kg}$.
\textbf{Đáp án của câu hỏi là D}
}
\end{ex}

\begin{ex} %[3D2V24]
Hiện nay mẹ 32 tuổi, con 4 tuổi. Hỏi hiện nay tuổi mẹ gấp mấy lần tuổi con?
\choice
{\True $8$ lần}
{$6$ lần}
{$7$ lần}
{$9$ lần}
\loigiai{
Tuổi mẹ gấp tuổi con số lần là: $32 : 4 = 8$ (lần).
\textbf{Đáp án của câu hỏi là A}
}
\end{ex}

\begin{ex} %[3D2V24]
Bạn Bình có 15 viên bi. Nếu Nam có thêm 3 viên bi thì số bi của Nam sẽ gấp 3 lần số bi của Bình. Hỏi Nam có bao nhiêu viên bi?
\choice
{$45$ viên bi}
{$48$ viên bi}
{\True $42$ viên bi}
{$39$ viên bi}
\loigiai{
Nếu Nam có thêm 3 viên bi thì Nam có: $15 \times 3 = 45$ viên. Số bi hiện tại của Nam: $45 - 3 = 42$ viên bi.
\textbf{Đáp án của câu hỏi là C}
}
\end{ex}

\textbf{Phần 2. Trắc nghiệm đúng sai}
\setcounterpageref{ex}{0}

\begin{ex} %[3D2V24]
Xét tính đúng sai của các câu nói về việc gấp một số lên nhiều lần:
\choiceTF
{\True Muốn gấp số $a$ lên $k$ lần ta tính $a \times k$.}
{Gấp số 15 lên 3 lần thì được kết quả là 18.}
{\True Gấp một số chẵn lên một số lần bất kì luôn được kết quả là số chẵn.}
{Gấp 1 lên 100 lần thì bằng 10.}
\loigiai{
\begin{itemize}
    \item Ý a) Đúng: Quy tắc cơ bản.
    \item Ý b) Sai: Sai vì $15 + 3 = 18$ là thêm 3 đơn vị, gấp 3 lần phải là $15 \times 3 = 45$.
    \item Ý c) Đúng: Số chẵn nhân số nào cũng ra số chẵn.
    \item Ý d) Sai: Sai vì $1 \times 100 = 100$.
\end{itemize}
\begin{itemize}
    \item a) \textbf{Đúng}.
    \item b) \textbf{Sai}.
    \item c) \textbf{Đúng}.
    \item d) \textbf{Sai}.
\end{itemize}
}
\end{ex}

\begin{ex} %[3D2V24]
Lớp 3B có 14 bạn nữ, số bạn nam gấp đôi số bạn nữ:
\choiceTF
{\True Lớp 3B có 28 bạn nam.}
{\True Tổng số học sinh lớp 3B là 42 bạn.}
{Số bạn nam nhiều hơn số bạn nữ là 28 bạn.}
{\True Số học sinh cả lớp gấp 3 lần số bạn nữ.}
\loigiai{
\begin{itemize}
    \item Ý a) Đúng: $14 \times 2 = 28$ bạn.
    \item Ý b) Đúng: $14 + 28 = 42$ bạn.
    \item Ý c) Sai: Số nam nhiều hơn số nữ là: $28 - 14 = 14$ bạn.
    \item Ý d) Đúng: $42 : 14 = 3$ lần.
\end{itemize}
\begin{itemize}
    \item a) \textbf{Đúng}.
    \item b) \textbf{Đúng}.
    \item c) \textbf{Sai}.
    \item d) \textbf{Đúng}.
\end{itemize}
}
\end{ex}

\textbf{Phần 3. Trả lời ngắn}
\setcounterpageref{ex}{0}

\begin{ex} %[3D2V24]
Gấp số 18 lên 4 lần rồi trừ đi 25 ta được số nào?
\shortans{47}
\loigiai{
\textbf{Đáp số:} $47$\\
$18 \times 4 = 72$; $72 - 25 = 47$.
}
\end{ex}

\begin{ex} %[3D2V24]
Bác Ba có 16 con vịt. Số con gà gấp 5 lần số con vịt. Hỏi bác Ba nuôi tất cả bao nhiêu con gà và vịt?
\shortans{96}
\loigiai{
\textbf{Đáp số:} $96$\\
Số gà: $16 \times 5 = 80$ con. Tổng gà và vịt: $16 + 80 = 96$ con.
}
\end{ex}

\begin{ex} %[3D2V24]
Một đoạn thẳng dài $12\text{ cm}$. Gấp đoạn thẳng đó lên 6 lần thì được đoạn thẳng dài bao nhiêu xăng-ti-mét?
\shortans{72}
\loigiai{
\textbf{Đáp số:} $72$\\
$12 \times 6 = 72\text{ (cm)}$.
}
\end{ex}

\begin{ex} %[3D2V24]
Năm nay con 6 tuổi. Tuổi mẹ gấp 5 lần tuổi con, tuổi bà gấp đôi tuổi mẹ. Hỏi năm nay bà bao nhiêu tuổi?
\shortans{60}
\loigiai{
\textbf{Đáp số:} $60$\\
Tuổi mẹ: $6 \times 5 = 30$ tuổi. Tuổi bà: $30 \times 2 = 60$ tuổi.
}
\end{ex}

\begin{ex} %[3D2V24]
Tìm một số biết rằng nếu gấp số đó lên 3 lần rồi bớt đi 12 thì được 45.
\shortans{19}
\loigiai{
\textbf{Đáp số:} $19$\\
Số đó nhân 3 bằng: $45 + 12 = 57$. Số cần tìm: $57 : 3 = 19$.
}
\end{ex}

\begin{ex} %[3D2V24]
Một hình chữ nhật có chiều rộng $15\text{ m}$, chiều dài gấp 3 lần chiều rộng. Tính chu vi của hình chữ nhật đó theo mét.
\shortans{120}
\loigiai{
\textbf{Đáp số:} $120$\\
Chiều dài: $15 \times 3 = 45\text{ m}$. Chu vi: $(45 + 15) \times 2 = 120\text{ m}$.
}
\end{ex}
